\section{Elicitation 2026.10.08}

\begin{tabular}{ll}
\textbf{Date:} & 2026.10.08 \\
\textbf{Consultant:} & Hector Hill / Gy'ax \\
\textbf{Elicitors:} & Zahra, Andrew, Matthias \\
\textbf{Location:} & TFS elicitation room \\
\textbf{Context:} & Third elicitation for LING 431/531 Field Methods \\
\textbf{Audio:} & \texttt{October\_08\_2026\_HH.WAV} \\
\end{tabular}

\subsection{Matthias [58:20-end]}

[58:48]: Henry glosses \textit{\tsb{g}alts'ep} as container-community since \textit{\tsb{g}al-} is a container-like prefix (e.g. grammar p. 62), but I will respectfully abstain. 

\begin{exe}
    \ex \glll hilthl Gitxsen la\tsb{x}/ts'im \tsb{g}alts'ep'm \\
    hilt=hl Gitxsen la\tsb{x}/ts'im \tsb{g}alts'ep-'m \\
    many=\Cn{} Gitksan on/in village-\Fpl{}.\Ii{}  \\
    `There's a lot of Gitksan in our village.'
\end{exe}

[1:00:32]:

\begin{exe}
    \ex \glll k'yep wilk'yephl Gitxsen la\tsb{x}/ts'im \tsb{g}alts'ep'm \\
    k'yep wil-k'yep=hl Gitxsen ts'im \tsb{g}alts'ep-'m \\
    ten times-ten=\Cn{} Gitksan in village-\Fpl{}.\Ii{}  \\
    `There's one hundred Gitksan in our village.'
\end{exe}

[1:01:08]: I think you need \textit{jo\tsb{g}it} here. But, without it wasn't overtly ruled out, so could try with a stronger context.

\begin{exe}
    \ex \glll hlibuuhl Gitxsen jo\tsb{g}it ts'im Vancouver \\
    hlibuu=hl Gitxsen jo\tsb{g}-t ts'im Vancouver \\
    few=\Cn{} Gitksan live-\Tsg{}.\Ii{} in Vancouver  \\
    `Few Gitksan live in Vancouver.'
\end{exe}

[1:02:50]: 

\begin{exe}
    \ex \glll \tsb{k}'am bagadilhl Gitxsen jo\tsb{g}it ts'im Vancouver \\
    \tsb{k}'am bagadil=hl Gitxsen jo\tsb{g}-t ts'im Vancouver \\
    only two=\Cn{} Gitksan live-\Tsg{}.\Ii{} in Vancouver  \\
    `Only two Gitksan live in Vancouver.'
\end{exe}

[1:03:37]: Hector accepted this, at the same time, I definitely did a bad job elicitng this, as I kept prodding without any real context. Take judgement with a grain of salt. Unquestionably preferred with \textit{jo\tsb{g}}.

\begin{exe}
    \ex \glll \tsb{k}'am bagadilhl Gitxsen ts'im Vancouver \\
    \tsb{k}'am bagadil=hl Gitxsen ts'im Vancouver \\
    only two=\Cn{} Gitksan  in Vancouver  \\
    `There are only two Gitksan in Vancouver.'
\end{exe}

[1:05:25]: Original prompt was without \textit{\tsb{g}a}, Hector responded with \textit{\tsb{g}a}. Don't know why exactly, but Rigsby p. 95 might have a clue.

Hector says \textit{\tsb{g}a} is ``specific, there may be brown bears but you are only counting the black bears''.

This was based on \textcite[314]{bicevskisQuantificationGitksan2017}, replacing \textit{hilt} with \textit{gilbil}. This environment is unambiguously a moved [Q NP] due to a ban on postnominal adjectives. Will have to think about a case where something is unambiguously predicative.

\begin{exe}
    \ex \glll hilthl smexhl \tsb{g}at'uuts'xwit \\
    hilt=hl smex=hl \tsb{g}a-t'uuts'xw-it \\
    many=\Cn{} bear=\Cn{} \Distr{}-black-\Sx{}   \\
    `There are many black bears.'
\end{exe}

[1:06:33]: It seems like \textit{\tsb{g}abiihl} is strongly preferred here, without it Hector just says ``it may work'', doesn't feel grammatical.

``If there's ten black bears, there has to be something else you put on there.\textit{ xbil smexhl \tsb{g}at'uuts'xwit} doesn't say anything.'' said at 1:09:45.

Also, it might be a clue that we do not hear \textit{=hl} on \textit{xbil} here. Could try putting it on to see if its accepted.

\begin{exe}
    \ex \glll xbil *(\tsb{g}abiihl) smexhl \tsb{g}at'uuts'xwit \\
    xbil \tsb{g}abii=hl smex=hl \tsb{g}a-t'uuts'xw-it \\
    ten  \Cnt{}=\Cn{} bear=\Cn{} \Distr{}-black-\Sx{}   \\
    `There are ten black bears.'
\end{exe}

[1:11:00]: \textbf{Context:} Telling Andrew the things I bought at a store. He already knows that I was buying shirts.

Hector repeating the construction at 1:12:53.

Hector says that this sentence sounds like I am translating from English... so there is no guarantee that this is actually grammatical. Instead, the sentence at 1:14:22 is offered.

\begin{exe}
    \ex \glll ?? anoogi'y sla'wxws dipun ii nagiikwhl hlibuut \\
    {} anoog-i-'y sla'wxws dip=un ii na=giikw=hl hlibuu-t  \\
    {} like-\Tr{}-\Fsg{}.\Ii{} shirt \Assoc{}=\Dem{}.\Prox{} \Ccnj{} \Fsg{}.\I{}=buy=\Cn{} few-\Tsg{}.\Ii{} \\
    `I liked these shirts and I bought a few.'
\end{exe}

[1:14:22]: This is interesting because we have a \textit{=hl} following the RC, but it'd be weird to think of the RC being the predicate? Or is it like `What I bought are the shirts that I like'?

\begin{exe}
    \ex 
    \begin{xlist}
    \item     \glll hlibuuhl sla'wxws anoogi'y=hl giikwi'y \\
    hlibuu=hl sla'wxws anoog-i-'y=hl giikw-i-'y \\
    few-\Cn{} shirt like-\Tr{}-\Fsg{}.\Ii{}=\Cn{} buy-\Tr{}-\Fsg{} \\
    \textit{HH offered: }`A few shirts that I like I bought.'

    \item     \glll hlibuuhl sla'wxws anoogi'y ii nagiikwit \\
    hlibuu=hl sla'wxws anoog-i-'y ii na=giikw-t  \\
    few-\Cn{} shirt like-\Tr{}-\Fsg{}.\Ii{}  \Ccnj{} \Fsg{}.\I{}=buy-\Tsg{}.\Ii{} \\
    `\textit{HH offered: }`I liked a few shirts that I bought.'
    \end{xlist}
\end{exe}

[1:16:56]: First one is ``okay, but not in the way we say it''.

\begin{exe}
    \ex 
    \begin{xlist}
    \item  \glll ?? giikwi'y hlibuut \\
    {} giikw-i-'y hlibuu-t \\
    {}  buy-\Tr{}-\Fsg{} few-\Tsg{}.\Ii{}  \\
    \textit{Intended: }`I bought a few.'

    \item     \glll hlibuuhl giigwi'y \\
    hlibuu=hl giikw-i-'y \\
    few=\Cn{} buy-\Tr{}-\Fsg{}.\Ii{} \\
    `I bought a few.'
    \end{xlist}
\end{exe}

[1:18:45]:

\begin{exe}
    \ex 
    \begin{xlist}
    \item  \glll * giigwi'y gilbilt \\
    {} giikw-i-'y gilbil-t \\
    {}  buy-\Tr{}-\Fsg{} two-\Tsg{}.\Ii{}  \\
    \textit{Intended: }`I bought two.'

    \item     \glll gilbilhl giigwi'y \\
    gilbil=hl giikw-i-'y \\
    two=\Cn{} buy-\Tr{}-\Fsg{}.\Ii{} \\
    `I bought two.'
    \end{xlist}
\end{exe}

[1:19:23]: \textbf{Context: }You are about to start a hike, and your friends ask you whether it is an easy hike.

Can use reduplicated \textit{hlephlebiksxw} to mean ``we will \textit{all} be tired'' here.

\begin{exe}
    \ex \glll dim hlebiksxw 'nuu'm \\
    dim hlebiksxw 'nuu'm \\
    \Prosp{} tired \Fpl{}.\Iii{} \\
    `We will be tired.'
\end{exe}

[1:20:39]: \textbf{Context: }Your friend invited your family to a party, and he asks who will come. You don't know how many family members will come, but you don't want to lowball it. 

I hear something like \textit{'nuu'm-s} with a clear \textit{s} after, but I have no clue what this \textit{s} would be? It could be \Val{}, but does it appear on pronouns?

\begin{exe}
    \ex \glll dim hilt 'nuu'm \\
    dim hilt 'nuu'm \\
    \Prosp{} many \Fpl{}.\Iii{} \\
    `We will be many.' \textit{Offered}: `There's gonna be a lot of us.'
\end{exe}

[1:22:05]:

\begin{exe}
    \ex 
    \begin{xlist}
        \item \glll ?? dim xwsdins 'nuu'm \\
     {} dim xwsdins \tsb{g}abii-'m \\
    {} \Prosp{} five \Fpl{}.\Iii{} \\
    \textit{Intended: }`We will be five.'

    \item \glll dim xwsdins \tsb{g}abii'm \\
     dim xwsdins \tsb{g}abii-'m \\
    \Prosp{} five \Cnt{}-\Fpl{}.\Ii{} \\
    `We will be five.'
    \end{xlist}
    
\end{exe}

[1:23:52]: \textbf{Context: } I am waiting in line to get my order from a coffee shop. Andrew is getting impatient, so he asks me “How many people are ahead of us?”. I try to count, but I don't know the exact number.

\begin{exe}
    \ex \glll k'aphl gabiihl gyet s\tsb{g}oo\tsb{g}at \\
    k'ap=hl gabii=hl gyet xs=\tsb{g}oo\tsb{g}-it  \\
    ten=\Cn{} \Cnt{}=\Cn{} person \Supr{}-front-\Sx{} \\
    `There are ten people in front.'
\end{exe}

[1:24:08]: I really don't know what that final syllable is. Hector says \textit{s\tsb{g}oo\tsb{g}atlam} is ``the one that's ahead/starting''.

\begin{exe}
    \ex \glll k'apimaa gabiihl gyet s\tsb{g}oo\tsb{g}at(lam???) \\
    k'ap=imaa gabii=hl gyet xs=\tsb{g}oo\tsb{g}-it  \\
    ten=\Epis{} \Cnt{}=\Cn{} person \Supr{}-front-\Sx{} \\
    `There are about ten people in front.'
\end{exe}
